\documentclass[12pt, a4paper]{article}

% ==================== 中文支持 ====================
\usepackage[UTF8]{ctex}

% ==================== 页面设置 ====================
\usepackage{geometry}
\geometry{left=2.5cm, right=2.5cm, top=2.54cm, bottom=2.54cm}

% ==================== 数学公式 ====================
\usepackage{amsmath, amssymb}

% ==================== 表格相关 ====================
\usepackage{booktabs}
\usepackage{multirow}
\usepackage{array}
\usepackage{tabularx}
% 定义居中且可自动换行的列类型
\newcolumntype{Y}{>{\centering\arraybackslash}X}

% ==================== 图片与浮动体 ====================
\usepackage{graphicx}
\usepackage{float}

% ==================== 代码与颜色 ====================
\usepackage{listings}
\usepackage{xcolor}

% ==================== 标题格式定制 ====================
\usepackage{titlesec}
% 一级标题居中
\titleformat{\section}{\centering\Large\bfseries}{\thesection}{1em}{}
% 二级标题左对齐
\titleformat{\subsection}{\flushleft\large\bfseries}{\thesubsection}{1em}{}
% 三级标题左对齐
\titleformat{\subsubsection}{\flushleft\normalsize\bfseries}{\thesubsubsection}{1em}{}



\begin{document}

% --- 论文题目 (居中) [cite: 1, 16] ---
\begin{center}
    \LARGE \bfseries 基于 XXX 模型的 XXX 问题研究
\end{center}

% ==================== 摘要 ====================
\section*{摘要}
\addcontentsline{toc}{section}{摘要}

（在此处填写摘要内容）

\vspace{1em}
\noindent \textbf{关键词：} 

\newpage

% ==================== 目录 ====================
\tableofcontents
\newpage

% ==================== 正文开始 ====================
% --- 一、问题重述 (居中) ---
\section{问题重述}
该部分主要是用自己的话对问题进行总结，\textbf{关键是改写} 。字数在半页到一页左右 。

% --- 二、问题分析 (居中)  ---
\section{问题分析}
问题分析起着承上启下的作用 。

\subsection{问题一的分析} % 二级标题顶格
分析过程要简明扼要，不需要放结论 。建议使用流程图列出思维过程 。

% --- 三、模型假设 (居中) ---
\section{模型假设}
\begin{itemize}
    \item 假设题目给出的测量数据准确，无异常值 。
    \item 整个系统温度恒定 。
\end{itemize}

% --- 四、符号说明 (居中)  ---
\section{符号说明}
一般以三线表形式给出，包括符号、含义和单位 。建议用希腊字母 。

\begin{table}[H]
    \centering
    \begin{tabular}{ccc}
        \toprule
        变量 & 说明 & 量纲 \\
        \midrule
        $t$ & 时间 & $s$  \\
        $a_{ij}$ & 在$(i,j)$位置的人数 & 人  \\
        \bottomrule
    \end{tabular}
\end{table}

% --- 五、模型建立与求解 (居中) ---
\section{模型建立与求解}
\subsection{数据预处理}
\subsection{描述性统计}
\subsection{问题一模型的建立与求解}
\subsubsection{模型选择与建立}
解释清楚数据特点是什么样的，为什么要选择这个模型，这个模型是解决什么问题的，
以及自己模型的推导过程
\subsubsection{模型的求解}
讲解你模型求解算法是什么，这个算法的优势是什么
\subsubsection{模型结果}



% --- 六、模型检验 (居中)  ---
\section{模型检验}
包括稳定性分析、敏感性分析及误差分析。

% --- 七、模型优缺点评价 (居中)  ---
\section{模型优缺点评价}
优点说充分，缺点不回避 。


% ==================== 参考文献与附录 ====================
\newpage
\begin{thebibliography}{99}
% 在此处添加参考文献条目
\end{thebibliography}

\newpage
\section*{附录}
\addcontentsline{toc}{section}{附录}
（在此处添加附录内容，如代码、表格等）

\end{document}